\section{Probability and Statistics}

\subsection{Probability via Area}

\begin{exercise}
If $x,y$ are randomly chosen in $[0, 1]$, what is the probability that
$x \geq 2y$?
\end{exercise}
\begin{solution}
This probability can be determined by a ratio of areas. The total space has area $1$.
The area of the condition $x \geq 2y$ is as follows:
\[
\int_{y = 0}^{\frac{1}{2}} \int_{x = 2y}^1 dx dy =
\int_{y = 0}^{\frac{1}{2}} 1 - 2y dy =
\frac{1}{2} - \left(\frac{1}{2}\right)^2 - \frac{1}{4}.
\]
Thus, the probability is $\frac{1}{4}$. Note that in this case, it's better to just 
draw it out in $\R^2$ and use the triangle area formula.
\end{solution}


\subsection{Basic Probability}

Let $\Omega$ be the sample space and let $A,B \subseteq \Omega$ be events.
Recall that the complement of $A$ is $A^c = \Omega \setminus A$, and we have 
\[
\P(A^c)=1-\P(A).
\]
We can also talk about unions and intersections of events. This gives us the 
principle of inclusion and exclusion (PIE):
\[
\P(A\cup B)=\P(A)+\P(B)-\P(A\cap B).
\]
\begin{remark}
If $A$ and $B$ are disjoint, then $\P(A\cap B)=0$, so the above reads
\[
\P(A\cup B)=\P(A)+\P(B).
\]
\end{remark}

\subsection{Conditional Probability}

Recall that we also had the notion of conditional probability.
If $\P(B)>0$, then
\[
\P(A\mid B)=
\frac{\Pp(A\cap B)}{\Pp(B)}.
\]
Equivalently,
\[
\Pp(A\cap B)= \Pp(A\mid B)\Pp(B).
\]
Symmetrically, if $\P(A) > 0$, then
\[
\Pp(A\cap B) = \Pp(B\mid A)\Pp(A).
\]
What is nice is that we can now define the notion of independence of events.
We say that events $A$ and $B$ are independent if
\[
\Pp(A\cap B)=\Pp(A)\Pp(B).
\]
In the case that $\P(B) > 0$, conditional probability tells us that independence
is equivalent to
\[
\P(A\mid B)=\P(A).
\]
\begin{remark}
Note that independence and disjointness are distinct.
If $A$ and $B$ are disjoint and both have positive probability, then
\[
\Pp(A\cap B)=0,
\]
while
\[
\Pp(A)\Pp(B)>0.
\]
\end{remark}

These lead to some very useful results:

\begin{theorem}[Law of Total Probability]
Suppose $B_1,\ldots,B_n$ form a partition of $\Omega$. Then
\[
\Pp(A)
=
\sum_{i=1}^n
\Pp(A\mid B_i)\Pp(B_i).
\]
\end{theorem}

\begin{theorem}[Bayes' Theorem]
For events $A$ and $B$ with positive probability,
\[
\Pp(A\mid B)
=
\frac{\Pp(B\mid A)\Pp(A)}
{\Pp(B)}.
\]
\end{theorem}

Combining the two previous, we have Bayes's generalized theorem:

\begin{theorem}
If $B_1,\ldots,B_n$ form a partition of $\Omega$, then
\[
\Pp(B_j\mid A)
=
\frac{\Pp(A\mid B_j)\Pp(B_j)}
{\sum_{i=1}^n \Pp(A\mid B_i)\Pp(B_i)}.
\]
\end{theorem}



\subsection{Random Variables}

\begin{df}[Discrete Random Variable]
A discrete random variable $X$ assigns a real number to each outcome in the
sample space. Its probability mass function (pmf) is
\[
p_X(x)=\Pp(X=x),
\]
where
\[
\sum_x p_X(x)=1.
\]
\end{df}

\begin{df}[Continuous Random Variable]
A continuous random variable $X$ has a probability density function (pdf) $f_X$
satisfying
\[
f_X(x)\geq0
\]
and
\[
\int_{-\infty}^{\infty}f_X(x) dx=1.
\]
For $a<b$,
\[
\Pp(a\leq X\leq b) = \int_a^b f_X(x) dx.
\]
\end{df}

\begin{remark}
For a continuous random variable,
\[
\Pp(X=a)=0.
\]
\end{remark}

We also have the notion of the cumulative distributive function.

\begin{df}[Cumulative Distribution Function]
The cumulative distribution function of $X$ is
\[
F_X(x)=\Pp(X\leq x).
\]
\end{df}

\begin{remark}
For a continuous random variable, this means
\[
F_X(x)
=
\int_{-\infty}^x f_X(t) dt,
\]
where $f_X$ is the pdf. Therefore, whenever $F_X$ is differentiable,
\[
f_X(x)=F_X'(x).
\]
\end{remark}

\subsubsection{Expected Values of Random Variables}

If $X$ is a discrete random variable,
\[
\boxed{\E[X] =\sum_x x\Pp(X=x)}.
\]
If $X$ is a continuous random variable,
\[
\boxed{\E[X] =\int_{-\infty}^{\infty}x f_X(x) dx}.
\]

\begin{remark}
More generally, for a function $g$,
\[
\boxed{\E[g(X)]
=
\sum_x g(x)\Pp(X=x)}
\]
in the discrete case, and
\[
\boxed{\E[g(X)]
=
\int_{-\infty}^{\infty}g(x)f_X(x)\,dx}
\]
in the continuous case.
\end{remark}

Note that expectation is linear (for both discrete and continuous cases). 
That is, for any (real) constants $a,b$,
\[
\boxed{\E[aX+bY] = a\E[X]+b\E[Y]}.
\]
Of course, more generally, this says
\[
\E\left[\sum_{i=1}^n a_iX_i\right]
=
\sum_{i=1}^na_i\E[X_i],
\]
where the $a_i$ are constants.

\begin{remark}
Note that independence is not required for the lienarity of expectation.
\end{remark}


\subsubsection{Variance}

There is also the notion of variance. If we denote $\mu=\E[X]$, then
\[
\boxed{\Var(X)=\E[(X-\mu)^2] =\E[X^2]-\E[X]^2}.
\]
The second equality is justified by linearity. With this, we can define the 
standard deviation:
\[
\boxed{\sqrt{\Var(X)}}.
\]
Clearly, for constants $a,b$,
\[
\boxed{\Var(aX+b)=a^2\Var(X)}.
\]
Particularly,
\[
\Var(X+b)=\Var(X).
\]
We can also define covariance of two random variables $X$ and $Y$:
\[
\boxed{\Cov(X,Y) = \E[(X-\E[X])(Y-\E[Y])] = \E[XY]-\E[X]\E[Y]}.
\]
For arbitrary $X$ and $Y$,
\[
\boxed{
\Var(X+Y)
=
\Var(X)+\Var(Y)+2\Cov(X,Y).
}
\]
If $X$ and $Y$ are independent, then we note that
\[
\Cov(X,Y)=0,
\]
so
\[
\boxed{
\Var(X+Y)
=
\Var(X)+\Var(Y).
}
\]
\begin{remark}
Note that $\Cov(X,Y) = 0$ does not in general imply that $X$ and $Y$ are independent.
\end{remark}


\subsection{Common Distributions}

\begin{example}[Indicator Random Variable]
For an event $A$, define
\[
I_A=
\begin{cases}
1, & \text{if $A$ occurs},\\
0, & \text{otherwise}.
\end{cases}
\]
Then
\[
\E[I_A]= 1 \cdot \P(A) + 0 \cdot \P(A^c) = \P(A).
\]
If
\[
X=I_1+\cdots+I_n
\]
counts how many of a collection of events occur, then linaerity implies
\[
\E[X]
=
\E[I_1]+\cdots+\E[I_n].
\]
\end{example}

\begin{example}
If
\[
X=
\begin{cases}
1, & \text{with probability $p$},\\
0, & \text{with probability $1-p$},
\end{cases}
\]
then
\[
\E[X]=p, \qquad\Var(X)=p(1-p).
\]
\end{example}

\begin{example}[Binomial Distribution]  
If $X$ counts the number of successes in $n$ independent trials, each with
success probability $p$, then $X\sim\Bin(n,p)$.  
Its pmf is
\[
\Pp(X=k) =\binom{n}{k}p^k(1-p)^{n-k}.
\]   
Moreover,
\[
\E[X]=np, \qquad \Var(X)=np(1-p).
\]  
\end{example}

\begin{example}[Geometric Distribution]
If $X$ is the number of trials required to obtain the first success, then
\[
\Pp(X=k)=(1-p)^{k-1}p, \qquad k=1,2,\ldots
\]
and
\[
\E[X]=\frac1p,
\qquad
\Var(X)=\frac{1-p}{p^2}.
\]
The geometric distribution has the interesting property
\[
\Pp(X>m+n\mid X>m) = \Pp(X>n).
\]
\end{example}

\begin{example}[Poisson Distribution]
If $X\sim\Pois(\lambda)$, then
\[
\Pp(X=k) = e^{-\lambda}\frac{\lambda^k}{k!}, \qquad
k=0,1,2,\ldots
\]
and
\[
\E[X]=\lambda,
\qquad
\Var(X)=\lambda.
\]
\end{example}

\begin{example}[Uniform Distribution]
If $X$ is uniformly distributed on $[a,b]$, then
\[
f_X(x)
=
\frac{1}{b-a},
\qquad
a\leq x\leq b.
\]
Moreover,
\[
\E[X]=\frac{a+b}{2}, \qquad \Var(X)
=
\frac{(b-a)^2}{12}.
\]
\end{example}